\documentclass[12pt,a4paper]{exam}
\usepackage[utf8]{inputenc}
\usepackage{amsmath,amssymb,amsfonts}
\usepackage{geometry}
\usepackage{enumitem}
\usepackage{graphicx}
\usepackage{hyperref}
\usepackage{xcolor}
\geometry{a4paper, top=2cm, bottom=2cm, left=2.5cm, right=2.5cm}
\hypersetup{colorlinks=true, linkcolor=blue, urlcolor=blue}
\renewcommand{\thequestion}{\textbf{\arabic{question}}}
\renewcommand{\questionlabel}{\thequestion.}
\pointname{ marks}
\pointsdroppedatright
\bracketedpoints
\setlength{\parindent}{0pt}
\setlength{\emergencystretch}{3em}
\allowdisplaybreaks
\newsavebox{\schmathbox}
\newcommand{\fitmath}[1]{%
  \sbox{\schmathbox}{$\displaystyle #1$}%
  \ifdim\wd\schmathbox>\linewidth
    \resizebox{\linewidth}{!}{\usebox{\schmathbox}}%
  \else
    \usebox{\schmathbox}%
  \fi}

\begin{document}

\thispagestyle{empty}
\begin{center}
  \textbf{\LARGE Diag} \\[0.3cm]
  \textbf{Time Allowed: 3 Hours} \hfill \textbf{Maximum Marks: 80} \\[0.3cm]
  \rule{\textwidth}{0.5pt}
\end{center}

\vspace{0.3cm}

\noindent\textbf{General Instructions:}
\begin{enumerate}
  \item {\normalfont\normalcolor This question paper contains 40 questions. All questions are compulsory.}
  \item {\normalfont\normalcolor Questions are grouped into sections A through E.}
  \item {\normalfont\normalcolor Use of calculators is NOT permitted.}
\end{enumerate}
\bigskip
\noindent\textbf{\large Section A: Multiple Choice \& Assertion-Reason (1 mark each)}

\begin{questions}
\question[1] The discriminant of the quadratic equation $2x^2 - 5x + 3 = 0$ is:
\begin{enumerate}[label=(\Alph*)]
  \item (A) 1
  \item (B) 2
  \item (C) 3
  \item (D) 4
\end{enumerate}
\vspace{0.15cm}

\question[1] Which of the following is a quadratic equation?
\begin{enumerate}[label=(\Alph*)]
  \item (A) $x^3 + 2x = 0$
  \item (B) $2x + 5 = 0$
  \item (C) $x^2 - 3x + 2 = 0$
  \item (D) $x^2 + \frac{1}{x} = 0$
\end{enumerate}
\vspace{0.15cm}

\question[4] Which class interval contains the median of the distribution?
\begin{enumerate}[label=(\Alph*)]
  \item (A) 140-160
  \item (B) 120-140
  \item (C) 160-180
  \item (D) 180-200
\end{enumerate}
\vspace{0.15cm}

\question[4] Suppose the mean consumption of the given data is 170 liters. Which of the following is likely true based on the relation between mean, median, and mode?
\begin{enumerate}[label=(\Alph*)]
  \item (A) The distribution is symmetric.
  \item (B) The distribution is positively skewed.
  \item (C) The distribution is negatively skewed.
  \item (D) The mode is greater than 170.
\end{enumerate}
\vspace{0.15cm}

\question[1] The decimal expansion of $\frac{147}{1200}$ will terminate after how many places of decimal?
\begin{enumerate}[label=(\Alph*)]
  \item (A) 2
  \item (B) 3
  \item (C) 4
  \item (D) 5
\end{enumerate}
\vspace{0.15cm}

\question[1] The discriminant of the quadratic equation $x^2 - 5x + 6 = 0$ is:
\begin{enumerate}[label=(\Alph*)]
  \item (A) 1
  \item (B) -1
  \item (C) 49
  \item (D) 0
\end{enumerate}
\vspace{0.15cm}

\question[1] The discriminant of the quadratic equation $2x^2 - 3x + 1 = 0$ is:
\begin{enumerate}[label=(\Alph*)]
  \item (A) 1
  \item (B) -1
  \item (C) 9
  \item (D) -7
\end{enumerate}
\vspace{0.15cm}

\question[1] The discriminant of the quadratic equation $2x^2 - 3x + 5 = 0$ is:
\begin{enumerate}[label=(\Alph*)]
  \item (A) $31$
  \item (B) $-31$
  \item (C) $41$
  \item (D) $-41$
\end{enumerate}
\vspace{0.15cm}

\question[1] The roots of the quadratic equation $x^2 - 7x + 10 = 0$ are:
\begin{enumerate}[label=(\Alph*)]
  \item (A) 2, 5
  \item (B) -2, -5
  \item (C) 2, -5
  \item (D) -2, 5
\end{enumerate}
\vspace{0.15cm}

\question[1] Which of the following is the discriminant of the quadratic equation $2x^2 - 4x + 3 = 0$?
\begin{enumerate}[label=(\Alph*)]
  \item (A) $\sqrt{40}$
  \item (B) $4$
  \item (C) $16 - 24$
  \item (D) $0$
\end{enumerate}
\vspace{0.15cm}

\question[1] Assertion (A): The decimal expansion of $\frac{7}{250}$ is terminating.
Reason (R): A rational number $\frac{p}{q}$ has a terminating decimal expansion if the prime factorisation of $q$ is of the form $2^n 5^m$, where $n$ and $m$ are non-negative integers.
\begin{enumerate}[label=(\Alph*)]
  \item (A) Both Assertion (A) and Reason (R) are true and Reason (R) is the correct explanation of Assertion (A)
  \item (B) Both Assertion (A) and Reason (R) are true but Reason (R) is NOT the correct explanation of Assertion (A)
  \item (C) Assertion (A) is true but Reason (R) is false
  \item (D) Assertion (A) is false but Reason (R) is true
\end{enumerate}
\vspace{0.15cm}

\question[1] Assertion (A): The HCF of two consecutive odd numbers is always 1.
Reason (R): Two consecutive odd numbers are co-prime.
\begin{enumerate}[label=(\Alph*)]
  \item (A) Both Assertion (A) and Reason (R) are true and Reason (R) is the correct explanation of Assertion (A)
  \item (B) Both Assertion (A) and Reason (R) are true but Reason (R) is NOT the correct explanation of Assertion (A)
  \item (C) Assertion (A) is true but Reason (R) is false
  \item (D) Assertion (A) is false but Reason (R) is true
\end{enumerate}
\vspace{0.15cm}

\question[1] Assertion (A): The quadratic equation $2x^2 - 5x + 3 = 0$ has two distinct real roots.
Reason (R): For a quadratic equation $ax^2 + bx + c = 0$, if $b^2 - 4ac > 0$, then the equation has two distinct real roots.
\begin{enumerate}[label=(\Alph*)]
  \item (A) Both Assertion (A) and Reason (R) are true and Reason (R) is the correct explanation of Assertion (A)
  \item (B) Both Assertion (A) and Reason (R) are true but Reason (R) is NOT the correct explanation of Assertion (A)
  \item (C) Assertion (A) is true but Reason (R) is false
  \item (D) Assertion (A) is false but Reason (R) is true
\end{enumerate}
\vspace{0.15cm}

\question[1] Assertion (A): The quadratic equation $x^2 - 3x + 5 = 0$ has no real roots.
Reason (R): For a quadratic equation $ax^2+bx+c=0$, if discriminant $D = b^2 - 4ac < 0$, then the equation has no real roots.
\begin{enumerate}[label=(\Alph*)]
  \item (A) Both Assertion (A) and Reason (R) are true and Reason (R) is the correct explanation of Assertion (A)
  \item (B) Both Assertion (A) and Reason (R) are true but Reason (R) is NOT the correct explanation of Assertion (A)
  \item (C) Assertion (A) is true but Reason (R) is false
  \item (D) Assertion (A) is false but Reason (R) is true
\end{enumerate}
\vspace{0.15cm}

\end{questions}
\bigskip
\noindent\textbf{\large Section B: Very Short Answer (2 marks each)}

\begin{questions}
\question[4] Calculate the median daily water consumption (in liters) using the formula.
\vspace{0.15cm}

\question[4] If the municipality decides that households consuming more than 200 liters per day need a special water‑saving advisory, how many households fall into that category?
\vspace{0.15cm}

\question[1] If $A = 2^3 \times 5^2$ and $B = 2^2 \times 5^3 \times 7$, find the HCF of $A$ and $B$.
\vspace{0.15cm}

\question[2] Divide the polynomial $p(x) = x^3 - 3x^2 + 5x - 3$ by $g(x) = x^2 - 2$ using the division algorithm. Find the quotient $q(x)$ and remainder $r(x)$.
\vspace{0.15cm}

\question[1] Find the discriminant of the quadratic equation $2x^2 - 4x + 3 = 0$.
\vspace{0.15cm}

\question[1] Find the discriminant of the quadratic equation $2x^2 - 3x + 1 = 0$.
\vspace{0.15cm}

\question[1] Find the discriminant of the quadratic equation $2x^2 - 4x + 3 = 0$.
\vspace{0.15cm}

\question[2] Find the discriminant of the quadratic equation $2x^2 - 4x + 3 = 0$. Hence, determine the nature of its roots.
\vspace{0.15cm}

\question[2] Find the nature of the roots of the quadratic equation $2x^2 - 4x + 3 = 0$.
\vspace{0.15cm}

\question[2] Find the nature of the roots of the quadratic equation $2x^2 - 4x + 3 = 0$.
\vspace{0.15cm}

\question[2] Find the discriminant of the quadratic equation $3x^2 - 5x + 2 = 0$ and hence determine the nature of its roots.
\vspace{0.15cm}

\question[2] Find the discriminant of the quadratic equation $3x^2 - 6x + 2 = 0$ and hence determine the nature of its roots.
\vspace{0.15cm}

\end{questions}
\bigskip
\noindent\textbf{\large Section C: Short Answer (3 marks each)}

\begin{questions}
\question[3] Solve the following pair of equations by reducing them to a pair of linear equations:
$\frac{2}{x} + \frac{3}{y} = 11$ and $\frac{5}{x} - \frac{4}{y} = -7$
\vspace{0.15cm}

\question[3] A train travels 360 km at a uniform speed. If the speed had been 5 km/h more, it would have taken 1 hour less for the same journey. Find the speed of the train.
\vspace{0.15cm}

\question[3] The speed of a boat in still water is 8 km/h. It takes the same time to travel 15 km upstream as it takes to travel 25 km downstream. Find the speed of the stream.
\vspace{0.15cm}

\question[3] The sum of the squares of two consecutive positive integers is 365. Find the integers.
\vspace{0.15cm}

\question[3] The sum of two numbers is 15. If the sum of their reciprocals is $\frac{3}{10}$, find the numbers.
\vspace{0.15cm}

\end{questions}
\bigskip
\noindent\textbf{\large Section D: Long Answer (4-5 marks each)}

\begin{questions}
\question[4] In a right triangle $ABC$, right-angled at $B$, $BD$ is perpendicular to $AC$. Prove that: (i) $AB^2 = AD \times AC$, (ii) $BC^2 = CD \times CA$, and (iii) $BD^2 = AD \times DC$.
\vspace{0.15cm}

\question[4] A motorboat whose speed in still water is $18$ km/h takes $1$ hour more to go $24$ km upstream than to return downstream to the same spot. Find the speed of the stream. Also, find the total time taken for the entire journey (upstream and downstream).
\vspace{0.15cm}

\question[4] If the roots of the quadratic equation $(a-b)x^2 + (b-c)x + (c-a) = 0$ are equal, prove that $2a = b + c$.
\vspace{0.15cm}

\question[4] Two water taps together can fill a tank in $9\frac{3}{8}$ hours. The tap of larger diameter takes 10 hours less than the smaller one to fill the tank separately. Find the time in which each tap can separately fill the tank.
\vspace{0.15cm}

\question[4] If the roots of the quadratic equation $(a - b)x^2 + (b - c)x + (c - a) = 0$ are equal, prove that $2a = b + c$.
\vspace{0.15cm}

\question[5] A train travels a distance of 480 km at a uniform speed. If the speed had been 8 km/h less, then it would have taken 3 hours more to cover the same distance. Find the usual speed of the train. (Note: Form a quadratic equation and solve it.)
\vspace{0.15cm}

\question[4] In $\triangle ABC$, $P$ and $Q$ are points on sides $AB$ and $AC$ respectively such that $PQ \parallel BC$. The area of $\triangle APQ$ is $16 \text{ cm}^2$ and the area of trapezium $\text{PBCQ}$ is $20 \text{ cm}^2$. Find the ratio $AP:PB$. If $AB = 12$ cm, find the lengths of $AP$ and $PB$.
\vspace{0.15cm}

\end{questions}
\bigskip

\end{document}
